%!TeX root = ../incremental_SS_Translation.tex

\chapter{Śārīrasthāna 7:  Distinctions between the colours of the vessels in the body}


\section{Introduction}
\texttt{Draft text by Jan Gerris.
    The Sanskrit text of this adhyāya is a lightly 
    edited transcript of NAK 5-333,
    waiting for full critical attention.}

\section{Literature} 

Meulenbeld offered an annotated overview of this chapter and a
bibliography of earlier scholarship to 2002 and, in his notes, citations
of the parallel passages in the \CS.\fvolcite{IA}[256--257]{meul-hist}   

The first draft of the following translation was prepared by Jan
Gerris.
% and later revised in collaboration Philipp Maas, John Nemec,
%Dominik Wujastyk, Madhusudan Rimal, Kenneth Zysk and Andrey Klebanov.

\newpage

\section{Translation}

\begin{translation}
    
    % First draft by Jan Gerris, [Sept 2026]
    
    \bigskip       
    
    \begin{tt}
        \raggedright
        \bigskip
        \marginpar{Draft tr.\ from here}
        
\item[1]

Now we shall expound the anatomical description and classification of the vessels.

\item[2]  vulgate only

\item[3]
There are seven hundred vessels.
By these, this body is moistened like a garden by water-conduits or a field by irrigation channels; through their contractions and expansions, they resemble the network of veins on the leaves of a tree, and they indeed originate from the navel, spreading upwards, downwards, and obliquely.

\item[4]
And on this subject, there is a verse:
As many vessels as there are in the bodies of living beings, they are all attached to the navel and spread out in all directions.

\item[5]   vulgate only

\item[6]
Of those, there are forty primary vessels; among them, ten carry Vāta, ten carry Pitta, ten carry Kapha, and ten carry blood.
Of those Vāta-carrying vessels situated in the seat of Vāta, there are one hundred and seventy-five.
There are an equal number of Pitta-carrying vessels in the seat of Pitta.
And there are the blood-carrying vessels in the liver and spleen.
And there are the Kapha-carrying vessels in the seat of Kapha.
In this manner, these vessels total seven hundred. 

\item[7]
Among them, there are twenty-five Vāta-carrying vessels in one leg.
By this, the other leg and the two arms are explained.
In the trunk there are thirty-four; of these, eight are located in the pelvis, supported by the anus and penis, two in each of the two sides, six in the back, an equal number in the abdomen, and ten in the chest.
There are forty-one above the clavicles; of those, eight are in the tongue.
Nine are in the nose.
Two are in the neck.
Four are in the two jaw-bones.
Six are in the eyes, three in the forehead, and four in the ears.
Two are in the temples.
Three are in the head; thus, these one hundred and seventy-five [vessels] of the wind-carrying vessels have been explained.
This same classification applies to the remaining [vessels]; thus, these seven hundred vessels have been explained.

\item[8]
And here are the verses on this subject:

\item[8.0]
When the wind is not vitiated and resorts to its own carrying vessels, then the strength, complexion, and vitality of the person are healthy, and the mind is likewise clear.

\item[8.1]
The wind, moving through its own vessels, produces non-obstruction of physical actions, clarity of the functions of the intellect, and other beneficial qualities.

\item[9]
However, when the provoked wind enters its own vessels, then various diseases born of Vāta arise in the person.

\item[9.add]
When, however, the bile is not aggravated and occupies its own carrying vessels, then the digestive fire remains unimpaired and properly digests food.

\item[10]
While moving through its own vessels, Pitta produces brilliance, an appetite for food, the kindling of the digestive fire, freedom from disease, and other beneficial qualities.

\item[11]
But when the aggravated Pitta enters those vessels, then various diseases of Pitta origin arise in the person.

\item[11.1]
When, however, the unaggravated Kapha enters its own vessels, the person's internal organs and joints function without disease.

\item[12]
Kapha, moving through its own vessels, produces moisture in the limbs, stability of the joints, strength, a lack of depression, and other beneficial qualities.

\item[13]
But when the aggravated Kapha enters those vessels, then various diseases of Kapha origin arise in the person.

[cf. 1938 ed. 3.7.13.1]
When, however, the blood is not aggravated and occupies its own carrying vessels, then the person correctly perceives the pleasant and unpleasant nature of touches.

\item[14]
Blood, while circulating through its own vessels, produces clarity of complexion, stability, and nourishment of the tissues, as well as other beneficial qualities.

\item[15]
But when the aggravated blood enters those vessels, then various diseases of blood origin arise in the person.

\item[15.1]
When unvitiated, they support the bodies with beneficial fluids; however, when vitiated, they lead to disorders—of this there is no doubt.

\item[16]
Indeed, there are no vessels that carry only Vata, nor any that carry only Pitta; all these vessels carry everything (all humours).

\item[17]
When the humours are corrupted, intensified, and flowing rapidly, their course is certainly directed upwards, downwards, and obliquely; therefore, all vessels are considered to carry everything. 

\item[18]
Among them, the vessels carrying Vāta are reddish-brown, those carrying Pitta are blue, those carrying blood are red, and those carrying Kapha are white.

\item[19]
Henceforth, I shall describe those vessels in humans that must not be pierced.

\item[20]  vulgate only

\item[21]
In the limbs, sixteen are unpierceable, and thirty-two in the trunk; above the collarbone, fifty are declared to be unpierceable.

\item[19cd]
From piercing them, deformity or even death will certainly occur quickly.

\item[22]
In that regard, there are one hundred vessels in a single leg; among them, one is Jaladhārā and three are internal; of those, two are named Urvī and one is named Lohitākṣa, and they are unpierceable; by this, the other leg and the two arms are explained.
There are thirty-two in the pelvis. 
Among those, eight are not to be operated upon.
Two in each of the two groin regions and two in each of the two hip-bone regions.
There are eight in each side; among those, one should avoid one upward-moving vessel on each side and two located at the joints on both sides.
There are twenty-four in the back, along both sides of the spinal column.
Among those, one should avoid the two upward-moving ones known as Bṛhatī.
There are the same number in the abdomen; among those, one should avoid the two on each side of the line of hair above the penis.
There are forty in the chest; among those, fourteen are not to be operated upon: two in the heart, and one should avoid the two located at the base of each breast.
One should avoid four on both sides of the breast.
Two each in the Apalāpa and Stambha regions; thus, there are thirty-two vessels unfit for surgery.
In the pelvis, back, abdomen, and chest, there are one hundred and sixty-four; in the head, there are twenty.
In the neck, among those, eight and four are designated as vital points.
Two in the two Kṛkāṭikās, the same number in the two Vidhuras; in the two jaws, eight each, but among those, one should avoid the two connecting the joints.
There are thirty-six in the tongue; among those, sixteen situated below are not to be operated upon, as they are the taste-carrying vessels.
There are twenty-four in the nose; among those, one should avoid the four adjacent to the nose, and among those same ones, one in the palate.
There are thirty-six in both eyes.
Among those, one should avoid the single vessel at each of the two outer corners of the eyes.
There are sixty in the forehead; among those, one should avoid the four that follow the hairline.
In the two whorls of hair, there is one vessel each; in the two shoulder blades, there is one each to be avoided; and in the two ears, there are five each.
Among those, one should avoid the single sound-carrying vessel in each of the two ears.
In the two temples, there are the same number; among those, one should avoid the single vessel located at each temple-joint.
There are twelve in the head; among those, one should avoid the two vessels in the two Utkṣepas.
One should avoid two in each of the hair-parting lines, and one in the Adhipati (crown of the head).

\item[22.add1]
And there are verses on this subject.
Among the seven hundred vessels in the bodies of embodied beings, one should specify ninety-eight as those not to be operated upon with surgical instruments.

\item[22.add2]
One should understand those that throb, those that carry fluid, those that are dependent on vital points, and those located in the ligaments and joints as vessels that must not be pierced.

\item[22.add3]
Deformity or even death will certainly result from the piercing of these.

\item[23]
Spreading out from the navel, the vessels pervade the entire body, just as the stalks and other parts of a lotus spread from its bulb through the water.

Thus concludes the seventh chapter in the [treatise] related to the body.

\end{tt}

\end{translation}